% capitulos/apendice.tex

\section{Dependências e Ambiente de Desenvolvimento Planejado}

O desenvolvimento do LibrasVision utilizará as seguintes bibliotecas Python, selecionadas com base nas soluções de maior desempenho identificadas na revisão bibliográfica:

\begin{table}[htb]
    \centering
    \caption{Dependências planejadas para o ambiente de desenvolvimento}
    \label{tab:dependencias}
    \begin{tabular}{|l|l|l|}
        \hline
        \textbf{Biblioteca} & \textbf{Versão mínima} & \textbf{Finalidade} \\
        \hline
        OpenCV (\texttt{opencv-python}) & 4.5+ & Captura e processamento de frames de vídeo \\
        MediaPipe & 0.10+ & Extração de landmarks corporais em tempo real \\
        TensorFlow / Keras & 2.10+ & Construção e treinamento da rede LSTM \\
        NumPy & 1.23+ & Manipulação de sequências de landmarks \\
        Scikit-learn & 1.1+ & Divisão treino/validação e métricas de avaliação \\
        Matplotlib & 3.6+ & Geração de gráficos de desempenho \\
        \hline
    \end{tabular}
\end{table}

\section{Arquitetura Planejada do Dataset}

O dataset a ser construído na etapa de coleta do TCC2 seguirá a estrutura de diretórios abaixo, separando os vídeos brutos por classe e os vetores de \textit{landmarks} extraídos para uso no treinamento:

\begin{verbatim}
dataset/
|-- sinais/
|   |-- biblioteca/
|   |   |-- usuario1_seq01.mp4
|   |   |-- usuario1_seq02.mp4
|   |   `-- ...
|   |-- prova/
|   |   |-- usuario1_seq01.mp4
|   |   `-- ...
|   `-- secretaria/
|       `-- ...
|-- landmarks/
|   |-- biblioteca.npy
|   |-- prova.npy
|   `-- secretaria.npy
`-- labels.csv
\end{verbatim}

Cada arquivo \texttt{.npy} conterá as sequências de 543 \textit{landmarks} tridimensionais extraídos pelo MediaPipe Holistic por frame, organizados em janelas de 30 frames por sinal gravado. O arquivo \texttt{labels.csv} associará cada sequência à sua classe e ao identificador do voluntário, permitindo a separação rigorosa entre conjuntos de treino e teste por sinalizador.

\section{Hiperparâmetros Iniciais Planejados do Modelo LSTM}

Os valores abaixo representam o ponto de partida para o treinamento do modelo LSTM, definidos com base nos parâmetros reportados por \citeonline{kamble:25} e \citeonline{varshini:25} em arquiteturas similares. Esses valores serão ajustados durante a fase experimental do TCC2 por meio de busca em grade (\textit{grid search}) ou ajuste manual orientado pelas curvas de aprendizado.

\begin{table}[htb]
    \centering
    \caption{Hiperparâmetros iniciais de referência para o modelo LSTM}
    \label{tab:hiperparametros}
    \begin{tabular}{|l|c|l|}
        \hline
        \textbf{Parâmetro} & \textbf{Valor inicial} & \textbf{Justificativa} \\
        \hline
        Unidades LSTM (1ª camada) & 64 & Capacidade moderada para sequências curtas \\
        Unidades LSTM (2ª camada) & 32 & Redução progressiva de dimensionalidade \\
        Taxa de aprendizado & 0,001 & Padrão Adam para tarefas de classificação \\
        Número de épocas & 100 & Com \textit{early stopping} por validação \\
        Tamanho do \textit{batch} & 32 & Equilíbrio entre velocidade e estabilidade \\
        Função de ativação & ReLU & Eficaz em camadas densas intermediárias \\
        Função de perda & Categorical Crossentropy & Adequada para classificação multiclasse \\
        Otimizador & Adam & Adaptativo, padrão em redes recorrentes \\
        Divisão treino/validação & 80\% / 20\% & Conforme protocolo da Seção 3.5 \\
        \hline
    \end{tabular}
\end{table}
